\documentclass[11pt, a4paper]{article}

\usepackage[utf8]{inputenc}
\usepackage[T1]{fontenc}
\usepackage[french]{babel}
\usepackage{lmodern}
\usepackage[left=2.5cm, right=2.5cm, top=3cm, bottom=3cm]{geometry}
\usepackage{amsmath, amssymb}

\title{\textbf{Résumé du Projet de Fin d'Études} \\ \large Approximation de Matrices de Rang Faible pour la Compression d'Images avec la DVS}
\author{Berrin KARAKAYA}
\date{Juin 2026}

\begin{document}

\maketitle

\vspace{1cm}

\begin{abstract}
    \noindent Dans ce projet, nous étudions comment les outils de la Théorème Spectral peuvent être utilisés pour analyser des matrices de grande dimension. L'objectif est de comprendre comment obtenir une approximation de rang faible d'une image numérique tout en minimisant la perte d'information.

    \begin{itemize}
        \item La motivation du travail sera présentée.
        \item Le cadre théorique sera introduit.
        \item Le théorème spectral sera présenté.
        \item La SVD sera expliquée.
        \item Le quotient de Rayleigh sera présenté.
        \item L’approximation de rang sera introduite.
        \item Le théorème d’Eckart-Young sera présenté.
        \item Une application de compression d’image sera montrée.
    \end{itemize}
\end{abstract}

\end{document}